\item \textbf{Выбор уставок устройства резервирования при отказе выключателя}

\begin{enumerate}[label=\arabic{section}.\arabic{enumi}.\arabic*, labelsep=4pt, leftmargin=0em, itemindent=65pt, parsep=3pt]
%\begin{enumerate}[label={}, labelsep=4pt, leftmargin=0pt, itemindent=57pt]

\item
Ток срабатывания УРОВ расчитывается по выражению:
\begin{equation*}
    I_{\textsl{с.з.}} = 0,1 \cdot I_{\textsl{ном,в}} = 0,1 \cdot 250 = 25\; (\textsl{А})  ,
\end{equation*}
\begin{eqexpl}[15mm]
    \item{$I_{\textsl{ном,в}}$} номинальный ток стороны ВН трансформатора, А;
\end{eqexpl}
С учетом рекомендаций пункта \ref{urov:punkt}.

Уставка задается в относительных единицах от номинального тока аналогового входа устройства, для этого полученное значение пересчитывается по следующему выражению:
\begin{equation*}
    I_{\textsl{ср}} = \frac{I_{\textsl{с.з.}}}{I_{\textsl{ном.вх.уст}}} \cdot \frac{I_{\textsl{ном.вт.ТТ}}}{I_{\textsl{ном.пер.ТТ}}} = \frac{25}{5} \cdot \frac{5}{400} = 0,06\; (\textsl{о.е.}),
\end{equation*}
\begin{eqexpl}[15mm]
    \item{$I_{\textsl{ном.вх.уст}}$} номинальный ток аналоговых входов устройства, 1 или 5 А;
    \item{$I_{\textsl{ном.вт.ТТ}}$} номинальный вторичный ток ТТ, А;
    \item{$I_{\textsl{ном.пер.ТТ}}$} номинальный первичный ток ТТ, А.
\end{eqexpl}

\item
Выдержка времени выбирается по выражению (\ref{eq:urov2}):
\begin{equation*}
t_{\textsl{УРОВ}} = t_{\textsl{откл.выкл.}}+t_{\textsl{воз.РЗ}}+t_{\textsl{ош.РВ}}+t_{\textsl{зап.}} = 0,06 + 0,025 + 0,025 + 0,1 = 0,21\; (\textsl{с}) ,
\end{equation*}

\item
Параметр \textbf{$I_{\textsl{ср}}$} <<Ток срабатывания>> функции УРОВ принимается \textbf{равным 0,06}.
\item
Параметр \textbf{$T_{\textsl{ср}}$} <<Выдержка времени срабатывания>> функции УРОВ принимается \textbf{равным 0,21 c}.

\end{enumerate}
